\section{Preliminary experimental evaluation of the Synthetic-RF method}
\label{sec:experimental-evaluation}

\subsection{Experimental Setup}

We conduct our experiments on the Cora dataset \cite{yang_revisiting_2016}, a citation network with 2 708 nodes, 5 429 edges and 1 433 binary features for each node. The task was inductive node classification. We compare our proposed Synthetic-RF method against random search. The methods were both implemented using the Optuna hyperparameter optimization framework \cite{akiba_optuna_2019}.

Each HPO method combination was evaluated independently 5 times to account for variability in performance due to random initialization and stochastic training processes. For both methods, we set a patience of 10 -- if no improvement in validation performance was observed for 10 consecutive trials, the optimization process was stopped early.

\subsection{Synthetic Dataset Generation}

The synthetic datasets used to pretrain the metamodel were generated as described in Section~\ref{sec:synthetic-rf}, taking Cora as the target dataset, so that the baseline statistics were \( \mathvec{s} = \left( 2708, \, 7, \, 3.90, \, 0.81, \, 1433 \right) \). We generated \( n_{\mathcal{D}} = 5 \) synthetic datasets, evaluating 30 hyperparameter configurations on each, which yields a meta-dataset \( \mathcal{A} \) of 150 observations. The configurations were drawn uniformly at random from the search space described below, without replacement, and the same 30 configurations were used for every synthetic dataset, so that the metamodel observes each of them across the whole range of generated dataset properties.

The perturbation strength was set to \( \varepsilon = 0.3 \), that is, each mimicked statistic was scaled by a factor drawn uniformly from \( \left[ 0.7, 1.3 \right] \). The node count was clipped to \( \left[ n_{\min}, n_{\max} \right] = \left[ 1500, 4000 \right] \), a range chosen to bracket Cora's 2 708 nodes, and the number of classes to at most \( k_{\max} = 20 \). The block sizes were drawn from a symmetric Dirichlet distribution with concentration \( \alpha = 5 \), which keeps the communities of comparable but not identical size. Since Cora's feature count is not perturbed, all synthetic datasets have 1 433 node features, matching the target; note, however, that these are real-valued, whereas Cora's are binary bag-of-words indicators.

\subsection{Search Spaces}

In our experiments, we used the well-studied and general framework of GraphSAGE \cite{hamilton_inductive_2017} as the base GNN architecture for node classification tasks. We defined a search space encompassing 8 hyperparameters that are critical to the performance of GraphSAGE models:
\begin{itemize}
	\item Number of layers: 1 to 3
	\item Hidden layer width: 16 to 64
	\item Aggregation function: mean, max
	\item Activation function: ReLU, PReLU
	\item Learning rate: \( 10^{-2} \), \( 5 \times 10^{-3} \), \( 10^{-3} \)
	\item Optimizer: Adam, AdamW
	\item \( L_2 \) regularization strength: 0, \( 5 \times 10^{-4} \), \( 5 \times 10^{-3} \)
	\item Dropout rate: 0, 0.5
\end{itemize}
The ranges were chosen based on the prior art \cite{hamilton_inductive_2017}. All experiments were trained for a maximum of 200 epochs, which was seldomly reached due to the fact that early stopping was employed to prevent overfitting.

\subsection{Metamodel Architecture}

Fot the metamodel, we used a scikit-learn Random Forest regressor with all parameter values taken as defaults from scikit-learn -- most importantly, 100 trees and 30\% of all features considered for each split. Several small implementation details are worth mentioning. First, in each step of the optimisation, all of the already evaluated configurations from \( \mathcal{A} \) are excluded from the set of candidate configurations \( \tilde{\Lambda} \), so that the same configuration is never needlessly re-evaluated. Second, all categorical features are one-hot encoded before being passed to the Random Forest model. Finally, in all our experiments, we used the F1-score as the performance metric \( \rho \) to be optimised.

\subsection{Results and Analysis}

The preliminary experimental evaluation is limited in nature. Since only one dataset and two methods were evaluated, the results are not sufficient to draw any general conclusions. However, the results do provide some initial insights into the performance of the Synthetic-RF method. On average (across 5 independent runs), Synthetic-RF achieved a final F1 score of 0.8628, while random search achieved an F1 score of 0.8535. This indicates that Synthetic-RF is able to find better hyperparameter configurations than random search, at least on the Cora dataset. However, the margin of improvement is modest at best, and further experiments on additional datasets and with more HPO methods are needed to validate the effectiveness of the Synthetic-RF method.

\subsection{Computational Cost}

The computational cost of Synthetic-RF splits into two parts: the trials on synthetic data that populate the meta-dataset \( \mathcal{A} \), and the trials on the target dataset performed during the optimisation itself. In the setup reported here, the synthetic datasets are generated to model the target dataset and are therefore a synthetic trial costs roughly as much as a trial on the target dataset. Including the pretraining, the method therefore consumes more computation than random search does.

The two parts differ in character, however. The pretraining trials are incurred once, ahead of time and independently of any particular tuning task, and they are mutually independent, so they may be spread across whatever hardware is available; a pretrained metamodel can then be reused for every subsequent tuning task, dividing its cost across all of them. The trials on the target dataset, in contrast, are what the practitioners care about: they are performed at tuning time and, for any sequential HPO method, they are inherently serial, each configuration being chosen based on those before it. It is this budget that pretraining on synthetic data is meant to reduce.
